\section{集成用平均与逐步加法改变误差}
\label{sec:13-Machine-Learning/03-Trees-Ensembles-and-Neighbors/02}

你已经知道一棵深树的方差有多大：训练数据微调几个样本，根节点的最佳切分就可能从``年龄<30''变成``收入<5000''，整棵树发生结构性改变。你试过剪枝、限制深度、加大最小叶样本数——这些都在降低方差的同时增加了偏差。

\textbf{困境是}：你在用只有一棵树的架构去同时处理方差和偏差，但它们属于不同来源的误差。有没有办法把方差的问题外包出去，而不用所有树必须是同一棵？

集成方法的思路是：既然一棵树不稳定，那就训练很多棵树。这些树不必每一棵都特别准，但只要它们的错误方向不同，取平均后剩下的误差就小于单棵。Bagging 和 Boosting 都基于这个直觉，但二者改变误差成分的方式截然不同。

\subsection{理解平均的方差削减：为什么独立才能获益}

设你有 \(M\) 个回归预测器，在某个 \(x\) 处的预测方差同为 \(\sigma^2\)，两两相关系数为 \(\rho\)。平均预测

\[
\bar f(x)=\frac1M\sum_{m=1}^{M}f_m(x)
\]

的方差为

\[
\Var(\bar f(x)) = \rho\sigma^2 + \frac{1-\rho}{M}\sigma^2.
\]

这一分解说明两件事：

第一，如果所有预测器完全相同（\(\rho=1\)），\(\Var(\bar f) = \sigma^2\)。复制同一个模型一万次，平均还是它自己——方差不变。

第二，即使 \(\rho>0\)，只要不到 1，随着 \(M\) 增大，第二项消失，残留的方差底限是 \(\rho\sigma^2\)。共同波动的部分永远不会被平均消除。

因此集成的核心矛盾不是``要多少棵树''，而是``怎样让树之间还不完全相关''。而要让预测器不同，要么扰动数据，要么扰动算法，要么扰动特征空间。

\subsection{Bagging：用 Bootstrap 制造不同的训练集}

Bagging 的做法简单直接：从原始训练集中有放回地抽取多个 bootstrap 样本，在每份样本上独立训练一棵树，然后对回归取平均、对分类投票。

一个 bootstrap 样本平均包含不同训练点的比例为

\[
1-\left(1-\frac1n\right)^n \longrightarrow 1-e^{-1}\approx 0.632.
\]

剩下约 36.8\% 的样本没有进入该棵树的训练集。这些样本可以做 out-of-bag（OOB）评价——对每个样本，只用那些没在训练集中见过它的树来投票或平均，近似得到交叉验证的误差估计。

但 bootstrap 的本质必须被认清：它不是在创造新的总体信息，而是在模拟训练数据的抽样波动。如果基学习器本身很稳定——比如你用的是线性回归而不是决策树——bootstrap 扰动后的预测几乎一模一样，那平均的收益就十分有限。Bagging 专门帮那些``数据微变就剧烈波动''的不稳定算法。

\subsection{随机森林：在 Bagging 之上再砍一刀相关性}

设想某个强特征——比如一个几乎完美预测的特殊指标——几乎所有 bootstrap 样本上，它都会在根节点被选中。于是各棵树的早期结构高度相似，树间相关 \(\rho\) 降不下来，方差削减碰到底限。

随机森林的做法是：在每个节点选择切分时，只允许\textbf{随机抽出的部分特征}参与候选。强特征不能在所有树的同一位置出现——有时它不在候选池里，其他特征就有机会上位。

这可能会略微增加单棵树的偏差（有时最优特征被暂时排除），但树间相关下降带来的方差削减通常远超过偏差增加。特征子采样数量 \(m_{\text{try}}\) 因此是一个关键参数：太小则单树太弱，太大则树间太像。

OOB 误差是一个便利的工具，但有两件事需要清醒认识：第一，如果特征工程和超参数选择已经看过了全部数据，OOB 就不再是真正独立的。第二，OOB 特征重要性——通过打乱某个特征后衡量 OOB 误差上升——会受到特征相关性的影响。两个特征可互相替代时，打乱其中一个，另一个可以顶上，边际损失很小；这不能推导出该特征没有预测信息。

\subsection{Boosting：换一个视角看集成}

Bagging 的成员可以并行训练，它们之间没有先后依赖。Boosting 完全不同：下一棵树专门修正前面所有树留下的残差。

考虑平方损失的加法模型：

\[
F_M(x) = \sum_{m=0}^{M} \eta\, h_m(x).
\]

在第 \(m\) 轮，当前模型 \(F_{m-1}\) 的残差是

\[
r_i^{(m)} = y_i - F_{m-1}(x_i).
\]

训练一棵小树 \(h_m\) 去预测这些残差，然后更新

\[
F_m(x) = F_{m-1}(x) + \eta\, h_m(x).
\]

残差是当前预测和目标之间的差值——让下一棵树去拟合这个差值，等价于沿平方损失的负梯度方向走一步。推广到一般可微损失时，残差换成负梯度：

\[
r_i^{(m)} = -\left.\frac{\partial \ell(y_i, F(x_i))}{\partial F(x_i)}\right|_{F=F_{m-1}}.
\]

这就是 gradient boosting：在函数值空间里沿经验风险的负梯度方向构造下一个基学习器。新树不必精确表示该梯度——它只是在允许的树类中做了一次近似拟合。重要的事情是，每一步的目标由损失函数的梯度规定，和 Bagging 的独立重采样原理完全不同。

\subsection{小步 + 浅树 = 控制函数增长速率}

学习率 \(\eta\)、树的深度和总轮数三个参数联合决定了 Boosting 的有效复杂度。较小的 \(\eta\) 通常需要更多轮才能收敛，但每次修正更温和——你不指望任何一棵树力挽狂澜，而是让许多棵树各自分担一小部分工作。深树允许单棵树捕捉高阶交互，但也更容易追逐噪声。

训练损失往往随轮数单调下降，但验证损失可能在某轮之后掉头上升——模型开始记忆训练集中的噪声结构。早停急在验证损失回升之前切断迭代。

但早停的时机必须由验证数据决定。先跑完所有轮数、画出验证曲线，再回头选最优停止位置——这是让测试集参与了训练过程。正确做法是在一个独立的验证折上实时监控，或至少把早停的决策纳入交叉验证的外层循环。

\subsection{Bagging 和 Boosting 处理的是不同成分的误差}

回到开头的困境：单棵树的方差大。Bagging 的核心操作是对不稳定的估计做平均，方差确实下降，偏差基本不变。Boosting 的核心操作是序贯的函数优化，每一步都在降低偏差——直到你走得太多，开始拟合噪声，方差反超。

把两者都概括为``多棵树投票''会丢掉最关键的差异：Bagging 希望树之间互不相同且各自偏差尚可，Boosting 希望每一步的修正温和且不跑过头。误差的降低公式、训练流程的并行性和对过拟合的防御手段，三者在两边都不相同。选择哪一个，取决于你的基学习器的问题是方差主导还是偏差主导。
